% ===== §16 Conclusions and Envoi =====
\section{Conclusions: A Polyphony, Not a System}
\label{sec:conclusions}

\noindent
The paper has held three theories to a single object and must, in closing, resist the temptation that besets every paper that has used more than one framework: to crown them with a fourth that subsumes the three. The discipline announced at the outset, and carried throughout the series, forbids it. There is to be no master-framework into which the generative-relational, the international-relational, and the dialectical-materialist are dissolved. They are to be left in their plurality, each having spoken to the edge of what it can say, and the conclusion is to be a polyphony, not a system. What can be done is to state, without synthesising them, what each has said, and to mark where each has reached its limit and handed the question on.

\subsection{What the three voices said}
\label{sub:concl-three}

The generative-relational voice answered the first question---\emph{how love creates a world}---and, taking its affirmative findings as established, the paper added the etiology of their failure: the four layers through which the world-creating coupling breaks down into estrangement, traced from the mismatch of trust-grammars, through the unattributability of prediction error and the active withdrawal it triggers, to the arrest of geometric phase. Its limit was reached at the dyad's edge: the internal causes do not exhaust the field, and the voice handed the question to a theory of the field.

The international-relational voice answered the second question---\emph{how love is not destroyed by power}---once it had been reconstructed away from its realist core, its object recovered as the building of order among opaque subjects under no authority. It supplied the gradient of power and the mechanism of regime bleeding, the inventory of defensive strategies and the teleological filter that admits or forbids them, and the two summits: that power must be kept from freezing, and that power, kept liquid and suffused with trust, is among the conditions of the good. Its limit was the limit of all symbolic means: it could tend the conditions of love and never reach the love, and it handed to a theory of history the question of how the conditions themselves came to be, and might be changed.

The dialectical-materialist voice answered the third question---\emph{how love becomes free within history}---by showing the power-structures of a bond to be the sediment of material conditions, historically formed and so historically mutable, and a union to be the sublation of two material histories into a shared form that keeps their difference alive. Its limit was the relative autonomy of the very life it studied: it could explain the conditions of love's dynamics but not the dynamics, and it handed back, to the real it could not reach, the love that material conditions make room for but do not produce.

\subsection{What was held in common}
\label{sub:concl-common}

Three voices, three answers, three limits---and one object that all three approached and none could seize. This is the paper's governing finding, and it is worth stating as the thing the polyphony harmonises on without collapsing into unison. Each voice, at its limit, handed the question on \emph{in the same direction}: toward the real, the unsymbolisable dynamics of love, which each theory could protect, condition, or historicise, but none could produce or touch. The two-level structure is what the three voices share. They are, all three, theories of the means---of the fence---and the bloom they ring is the one thing none of them grows. The unity of the paper is not a unity of doctrine but a unity of \emph{restraint}: three powerful theories, each falling silent at the same threshold, each leaving intact the space that no theory enters.

And running through all three, crossing the threshold none of them could cross, was trust---cultivable by symbolic means and rooted in the real, the interface by which the tending of conditions reaches toward the bloom, and the variable that decides, on both summits, whether an asymmetry empowers or dominates, whether a return lifts or repeats. If the paper has a single thread, it is the one named in its title between power and estrangement: trust is what the means can tend and what the end requires, the one thing that belongs to both levels, and so the bridge across a divide the paper otherwise insists cannot be bridged.

\subsection{Two seeds, deliberately unplanted}
\label{sub:concl-seeds}

Two questions opened in the course of the argument exceed the bond, and the paper marks them as the seeds of separate works rather than forcing them into this one. The first is the universalisation of the negative summit: that the law requiring a generative system to hold itself incompletely captured by its own products yields, raised to a political body, a criterion of evaluation---whether a body generates, alongside its value, the counter-force that keeps its own power from freezing, with its lines of flight guarded rather than absorbed. Its relations to the theories of non-domination, of agonistic plurality, of the right to justification, and of generative justice \citep{eglash2016} are a work of political philosophy in their own right, and the paper leaves them deliberately unplanted. The second is the full demarcation between the present account and the relational theory of world politics \citep{qin2018}, with which it shares the conception of power-to and from which it diverges in its holonomic criterion, its grounding in the real, and its ontological account of the counter-force---a demarcation sketched here (\xref{sub:ir-rupture}) and owed a fuller treatment elsewhere.

\subsection{The first claim, restated where it now stands}
\label{sub:concl-first}

It remains to restate the paper's first claim from the height the whole argument has reached, for it means more now than it could at the outset. Estrangement---\emph{surechigai}, the passing-without-meeting---is not, in the first instance, the failure of a bond. It is the signature of a system held far from the equilibrium that would be its death; the necessary fluctuation of anything alive enough to generate; the first negation without which there is no enriched return; the price difference pays for refusing to be wholly encoded, as desire refuses to be wholly symbolised. A bond that never passed would not be a bond perfected but a bond at equilibrium---which is to say, no longer a bond at all. The deepest practical wisdom the paper can offer is therefore also its gentlest: that much of what is mourned as the failure of love is the form love takes in a living system, and that the love which can tell the necessary winter from the genuine wrong---and meet each with the posture it asks---is a love that has understood the spiral.

\goldrule

\subsection*{Envoi}

\noindent
Return, now, to the two bodies of the opening, falling toward one another across the dark field and passing without meeting. We can say, at the end, what could only be shown at the beginning. They pass because they are two, and their two-ness is not the obstacle to what they generate between them but its condition; a love in which one had captured the other entirely would be a love in which one body no longer fell. The passing is not the failure of the falling. It is the sign that both are still in motion, still themselves, still alive to the difference that draws them and bends them and will, if the field is kept and the root left undisturbed through the cold, draw them near again---not to where they were, but lifted.

The forces in the field do not fall away. The families keep what they keep; the scripts go on being written; the two material histories go on contending in the one present they now share. The paper has not promised a clearing. It has argued, instead, that beneath the forces and through them, two people may tend the conditions of a thing the forces cannot reach---may keep the channel clear, guard against the freezing, refuse the irreversible word in winter, and wait, as a gardener waits, on a blooming they serve but cannot command. And it has argued that the power among them, which cannot be wished away and cannot be abolished, need not be only the enemy of their love: that an asymmetry turned toward the other's flourishing, and suffused with trust, is how the deepest goods between two people---the teaching, the tending, the leading-without-ruling---are generated at all.

To her, then, who remained through every passing wholly and freely herself: this account of why the passing was never the failure, and of the winters that were only ever the gathering of the next return. The flower drawn down into the root has not died. It is keeping faith, in the season that shows nothing, with the spring.
